\chapter{Business German: Geschäftskommunikation}
\label{cha:Business}

Dieses Kapitel behandelt die formelle Geschäftskommunikation auf Deutsch.

\section{Geschäftsbriefe (Business Letters)}
\label{sec:Geschaeftsbriefe}

\bigskip
\noindent\textbf{Struktur eines Geschäftsbriefes:}

\begin{center}
\begin{tabular}{|l|p{8cm}|}
\hline
Element & Beschreibung \\ \hline
Absender & Firmenname, Logo, Anschrift \\ \hline
Datum & Ort und Datum (rechtsbündig) \\ \hline
Empfänger & Name, Titel, Firma, Anschrift (linksbündig) \\ \hline
Betreff & Präzise und aussagekräftig \\ \hline
Anrede & Sehr geehrte/r... (formell) \\ \hline
Fließtext & Kerninhalt, strukturiert \\ \hline
Schlussformel & Mit freundlichen Grüßen / Beste Grüße \\ \hline
Unterschrift & Name in Druckschrift darunter \\ \hline
Anhänge & ggf. Vermerk „Anlage" \\ \hline
\end{tabular}
\end{center}

\bigskip
\noindent\textbf{Anrede-Formeln:}

\begin{itemize}
    \item \textbf{Sehr geehrte Damen und Herren,} (unbekannter Empfänger)
    \item \textbf{Sehr geehrter Herr Müller,} (bekannt, männlich)
    \item \textbf{Sehr geehrte Frau Schmidt,} (bekannt, weiblich)
    \item \textbf{Sehr geehrte Frau Dr. Weber,} (mit Titel)
    \item \textbf{Liebe Kolleginnen und Kollegen,} (informeller)
    \item \textbf{Sehr geehrte Geschäftspartner,} (bei Geschäftsbeziehungen)
\end{itemize}

\bigskip
\noindent\textbf{Schlussformeln:}

\begin{itemize}
    \item \textbf{Mit freundlichen Grüßen} (standard)
    \item \textbf{Mit freundlichen kollegialen Grüßen} (kollegial)
    \item \textbf{Beste Grüße} (etwas weniger formell)
    \item \textbf{Hochachtungsvoll} (sehr formell, selten)
    \item \textbf{Viele Grüße} (informell, nur bei bekannten Personen)
\end{itemize}

\section{E-Mail-Kommunikation}
\label{sec:Email}

\bigskip
\noindent\textbf{Betreffzeilen (Subject):}

\begin{center}
\begin{tabular}{|l|p{8cm}|}
\hline
Situation & Betreff \\ \hline
Anfrage & Anfrage bezüglich... / Anfrage: ... \\ \hline
Antwort & AW: Ihre Anfrage vom... / Re: ... \\ \hline
Auftrag & Auftrag: Lieferung von... \\ \hline
Termin & Terminbestätigung / Einladung zur Besprechung \\ \hline
Erinnerung & Erinnerung: Zahlung fällig \\ \hline
\end{tabular}
\end{center}

\bigskip
\noindent\textbf{Nützliche E-Mail-Phrasen:}

\begin{itemize}
    \item \textbf{Einleitung}: Hiermit... / Bezugnehmend auf... / Vielen Dank für Ihre E-Mail vom...
    \item \textbf{Überleitung}: bezüglich... / in Bezug auf... / hinsichtlich...
    \item \textbf{Handlungsaufforderung}: Wir bitten Sie um... / Bitte teilen Sie uns mit... / Wir würden uns freuen, wenn...
    \item \textbf{Schluss}: Bei Fragen stehe ich Ihnen gerne zur Verfügung. / Ich freue mich auf Ihre Rückmeldung.
\end{itemize}

\section{Besprechungen und Meetings}
\label{sec:Besprechungen}

\bigskip
\noindent\textbf{Phrasen für Meetings:}

\begin{center}
\begin{tabular}{|l|p{8cm}|}
\hline
Phase & Phrasen \\ \hline
Eröffnung & Die Sitzung wird eröffnet. / Wir beginnen mit... \\ \hline
Tagesordnung & Zunächst möchte ich auf die Tagesordnung eingehen. \\ \hline
Thema einbringen & Ich möchte das Thema... ansprechen. \\ \hline
Meinung äußern & Meiner Meinung nach... / Ich vertrete die Ansicht, dass... \\ \hline
Zustimmen & Dem kann ich nur zustimmen. / Das sehe ich genauso. \\ \hline
Widersprechen & Ich sehe das anders. / Dem muss ich widersprechen. \\ \hline
Nachfragen & Könnten Sie das bitte näher erläutern? \\ \hline
Zusammenfassung & Zusammenfassend lässt sich sagen... \\ \hline
Abschluss & Damit ist die Besprechung beendet. \\ \hline
\end{tabular}
\end{center}

\bigskip
\noindent\textbf{Typische Meeting-Vokabular:}

\begin{itemize}
    \item \textbf{Tagesordnung} - Agenda
    \item \textbf{Protokoll} - Minutes
    \item \textbf{Beschluss} - Decision
    \item \textbf{Maßnahme} - Measure/Action
    \item \textbf{Frist} - Deadline
    \item \textbf{To-do-Liste} - Task list
    \item \textbf{Rückmeldung} - Feedback
    \item \textbf{Vorschlag} - Proposal
    \item \textbf{Konsens} - Consensus
\end{itemize}

\section{Telefonate im Geschäftsumfeld}
\label{sec:Telefon}

\bigskip
\noindent\textbf{Telefonformulierungen:}

\begin{itemize}
    \item \textbf{Anruf annehmen}: Guten Tag, Firma X, mein Name ist...
    \item \textbf{ Sich melden}: ...spricht. / Hier ist...
    \item \textbf{Verbinden}: Ich verbinde Sie mit... / Einen Moment, bitte.
    \item \textbf{Um Rückruf bitten}: Kann Herr/Frau ... Sie zurückrufen?
    \item \textbf{Nachricht aufnehmen}: Möchten Sie eine Nachricht hinterlassen?
    \item \textbf{Gespräch beenden}: Vielen Dank für Ihren Anruf. / Auf Wiederhören.
\end{itemize}

\section{Bewerbung und Lebenslauf}
\label{sec:Bewerbung}

\bigskip
\noindent\textbf{Bewerbungsunterlagen:}

\begin{center}
\begin{tabular}{|l|p{8cm}|}
\hline
Dokument & Inhalt \\ \hline
 Anschreiben & \begin{minipage}{8cm}Motivation, Bezug zur Stelle, Kernqualifikationen\end{minipage} \\ \hline
 Lebenslauf & \begin{minipage}{8cm}Persönliche Daten, Berufserfahrung, Ausbildung, Kenntnisse\end{minipage} \\ \hline
 Zeugnisse & \begin{minipage}{8cm}Arbeitszeugnisse, Abschlusszeugnisse\end{minipage} \\ \hline
\end{tabular}
\end{center}

\bigskip
\noindent\textbf{Nützliche Bewerbungsphrasen:}

\begin{itemize}
    \item \textbf{Bezug auf Stellenanzeige}: Bezugnehmend auf Ihre Anzeige in... vom...
    \item \textbf{Motivation}: Ich bewerbe mich, weil... / Ich reizt die Herausforderung...
    \item \textbf{Qualifikationen}: Über mehrjährige Erfahrung in... verfüge ich. / Meine Stärken liegen in...
    \item \textbf{Schluss}: Über eine Einladung zu einem persönlichen Gespräch würde ich mich freuen.
\end{itemize}

\chapterend

\chapter{Academic German: Wissenschaftssprache}
\label{cha:Academic}

\section{Wissenschaftliche Texte}
\label{sec:WissenschaftlicheTexte}

\bigskip
\noindent\textbf{Typische Textsorten:}

\begin{center}
\begin{tabular}{|l|p{8cm}|}
\hline
Textart & Merkmale \\ \hline
Abstract & \begin{minipage}{8cm}Zusammenfassung, 150-300 Wörter\end{minipage} \\ \hline
Hausarbeit & \begin{minipage}{8cm}Einleitung, Hauptteil, Schluss, Literaturverzeichnis\end{minipage} \\ \hline
Bachelor/Masterarbeit & \begin{minipage}{8cm}Forschungsfrage, Methodik, Analyse, Fazit\end{minipage} \\ \hline
Essay & \begin{minipage}{8cm}Kurze Abhandlung, eigene Argumentation\end{minipage} \\ \hline
\end{tabular}
\end{center}

\section{Argumentation und Formulierungen}
\label{sec:Argumentation}

\bigskip
\noindent\textbf{Wissenschaftliche Wendungen:}

\begin{center}
\begin{tabular}{|l|p{8cm}|}
\hline
Funktion & Wendung \\ \hline
These aufstellen & Es wird argumentiert, dass... / Die These lautet... \\ \hline
Beispiel anführen & Beispielsweise... / Zum Beispiel... / Ein Beispiel hierfür ist... \\ \hline
Zitat & Laut... / Wie X (2020) schreibt... / „..." (X 2020: 55) \\ \hline
Ergebnis zusammenfassen & Zusammenfassend lässt sich feststellen... / Die Ergebnisse zeigen... \\ \hline
Eigene Meinung & Meines Erachtens... / Ich vertrete die Auffassung... \\ \hline
Einschränkung & Allerdings muss einschränkend gesagt werden... \\ \hline
\end{tabular}
\end{center}

\bigskip
\noindent\textbf{Übergangsphrasen:}

\begin{itemize}
    \item \textbf{Hinzufügen}: Des Weiteren... / Ferner... / Außerdem...
    \item \textbf{Gegensatz}: Allerdings... / Jedoch... / Demgegenüber...
    \item \textbf{Ursache}: Aufgrund... / Wegen... / Aus diesem Grund...
    \item \textbf{Folge}: Dadurch... / Infolgedessen... / Folglich...
    \item \textbf{Zeitlich}: Zunächst... / Anschließend... / Abschließend...
\end{itemize}

\section{Zitieren und Quellenangaben}
\label{sec:Zitieren}

\bigskip
\noindent\textbf{Zitierstile:}

\begin{itemize}
    \item \textbf{Harvard}: (Müller 2020, S. 55)
    \item \textbf{APA}: (Müller, 2020)
    \item \textbf{ISO 690}: Müller, Hans (2020): Titel. Ort: Verlag.
\end{itemize}

\bigskip
\noindent\textbf{Quellenarten:}

\begin{itemize}
    \item \textbf{Buch}: Autor, Vorname (Jahr): Titel. Untertitel. Ort: Verlag.
    \item \textbf{Artikel}: Autor, Vorname (Jahr): Titel. In: Zeitschrift, Band(Hef t), Seiten.
    \item \textbf{Online}: Autor (Datum): Titel. Online verfügbar unter: URL.
\end{itemize}

\section{Präsentationen}
\label{sec:Praesentationen}

\bigskip
\noindent\textbf{Präsentationsphrasen:}

\begin{center}
\begin{tabular}{|l|p{8cm}|}
\hline
Phase & Phrasen \\ \hline
Einleitung & \begin{minipage}{8cm}Mein Thema heute ist... / Ich möchte Ihnen zeigen...\end{minipage} \\ \hline
Gliederung & \begin{minipage}{8cm}Mein Vortrag gliedert sich in drei Teile...\end{minipage} \\ \hline
Übergang & \begin{minipage}{8cm}Nun möchte ich zum nächsten Punkt übergehen... / Damit komme ich zu...\end{minipage} \\ \hline
Illustration & \begin{minipage}{8cm}Wie Sie hier sehen... / Dies verdeutlicht...\end{minipage} \\ \hline
Schluss & \begin{minipage}{8cm}Zusammenfassend kann ich sagen... / Vielen Dank für Ihre Aufmerksamkeit.\end{minipage} \\ \hline
Fragen & \begin{minipage}{8cm}Ich freue mich auf Ihre Fragen. / Haben Sie dazu Fragen?\end{minipage} \\ \hline
\end{tabular}
\end{center}

\chapterend